\documentclass{article}
\usepackage[T1]{fontenc}
\usepackage{lmodern}
\usepackage{graphicx}
\usepackage{amsmath,amssymb}
\usepackage[a4paper,margin=2cm]{geometry}
\usepackage[absolute,overlay]{textpos}
\setlength{\TPHorizModule}{1mm}
\setlength{\TPVertModule}{1mm}
\usepackage[indonesian]{babel}
\usepackage{hyperref}
\setlength{\emergencystretch}{2em}
% unit: Halaman 16

\begin{document}

% === Halaman 16 (llm) ===

\textbf{Contoh 2:} Alice membangkitkan pasangan kuncinya. Bob akan mengirim pesan dengan menggunakan kunci publik Alice.

\textbf{(a) Pembangkitan kunci (oleh Alice)}

Alice memilih bilangan prima $p = 2273$, akar primitif dari $p$ yaitu $g = 3$, dan $x = 243$.

Alice kemudian menghitung:
\[ y = g^x \pmod{p} = 3^{243} \pmod{2273} = 461 \]

Jadi,

\begin{itemize}
    \item kunci privat Alice: $(x = 243, p = 2273)$
    \item kunci publik Alice: $(y = 461, g = 3, p = 2273)$.
\end{itemize}

\end{document}
